\documentclass[a4paper, 11pt]{article}
\usepackage[margin=0.75in]{geometry}
\usepackage{graphicx, tabularx, hyperref, amsmath, wrapfig}
\usepackage[usenames,dvipsnames]{color,xcolor}
\usepackage{listings}
\usepackage{hyperref}
\usepackage[siunitx]{circuitikz}
\usepackage{booktabs}
\usepackage{pgfplotstable, pgfplots}
\renewcommand{\arraystretch}{1.2}
\usepackage{xepersian}
\settextfont[Scale=0.9]{Dibaj}
\setdigitfont[Scale=1]{Dibaj}

\lstset{
    basicstyle=\ttfamily\small,
    keywordstyle=\color{blue},
    commentstyle=\color{green!50!black},
    stringstyle=\color{red!70!black},
    numbers=left,
    numberstyle=\tiny,
    stepnumber=1,
    breaklines=true,
    frame=single
}


\begin{document}

\begin{center}
\textbf{{\huge گزارش‌کار آزمایش شمارهٔ
4:
باردار شدن و بی‌بار شدن خازن‌ها
}}

\end{center}

\hrule
\noindent\begin{tabular}{rrl}
اعضای گروه: & ابوالفضل احمدی & 403172283\\
& نام و نام‌خانوادگی & 403001122
\end{tabular}\hfill
\begin{tabular}{rr}
تاریخ انجام آزمایش: & 24/08/1404 \\
دستیار آموزشی: &  \\
\end{tabular}\hfill
\begin{tabular}{rr}
گروه:  &  \\
زیرگروه: \lr{A} &  \\
\end{tabular}
\hrule

\section*{اهداف و خلاصهٔ شرح آزمایش}

\textbf{هدف‌ها:}
\begin{itemize}
  \item بررسی رفتار زمانیِ بارگیری و تخلیهٔ خازن‌ها و تأیید رابطهٔ نمایی \(V(t)=V_0 e^{-t/\tau}\).
  \item استخراج ثابت زمانیِ \(\tau=RC\) با روشِ کمینهٔ مربعات از داده‌های تجربی.
  \item برآورد مقاومت داخلیِ ولت‌متر \(R_\mathrm{v}\) در دو آزمایش مستقل و محاسبهٔ میانگین آن.
  \item تعیین ظرفیتِ معادلِ خازن‌ها در اتصالِ سری و موازی و مقایسه با روابط تئوری.
  \item محاسبهٔ خطای نسبیِ نتایجِ تجربی نسبت به مقادیر تئوری و بحث دربارهٔ منابعِ خطا.
\end{itemize}

\textbf{خلاصه:}
در این گزارش، با ثبت تغییراتِ ولتاژِ اندازه‌گیری‌شده توسطِ ولت‌متر در گذرِ زمان، برای دو پیکربندی متفاوت (باردار کردن خازن و تخلیهٔ خازن)، از برازشِ خطیِ \(\log_{10}(V/V_0)\) بر حسبِ \(t\) استفاده شد و شیبِ به‌دست‌آمده برای تعیینِ \(\tau\) و سپس \(R_\mathrm{v}\) به‌کار رفت. برای خازن \(C_1 = 20~\mu\mathrm{F}\) مقدارِ \(R_\mathrm{v} \approx 9~\mathrm{M\Omega}\) و برای خازن \(C_2 = 4~\mu\mathrm{F}\) مقدارِ \(R_\mathrm{v} \approx 12~\mathrm{M\Omega}\) به‌دست آمد و میانگین \(\bar{R} \approx 10.5~\mathrm{M\Omega}\) محاسبه شد. سپس ظرفیتِ معادل در اتصالِ سری و موازیِ \(C_1\) و \(C_2\) با استفاده از \(\tau\) تجربی و \(\bar{R}\) محاسبه و با روابطِ تئوریِ \(\tfrac{1}{C}=\tfrac{1}{C_1}+\tfrac{1}{C_2}\) (سری) و \(C=C_1+C_2\) (موازی) مقایسه شد. اختلاف‌ها حدود \(9\%-11\%\) بودند که با محدودیتِ دقتِ قرائتِ ولت‌متر، وابستگیِ مقاومتِ ورودی به محدودهٔ اندازه‌گیری، مقاومت‌های ناچیزِ سیم‌ها و اتصالات، و تقریب‌های عددی در برازش توجیه‌پذیر است. در مجموع، نتایج با مدل تئوری سازگارند و روشِ تعیینِ \(R_\mathrm{v}\) و \(C_\mathrm{eq}\) از رویِ \(\tau\) کارآمد و قابلِ اتکا ارزیابی می‌شود.

\section{باردار شدن خازن $C_1$ و تعیین مقاومت داخلی ولت‌متر }
\begin{wraptable}{l}{4cm} % r = راست جدول، 4cm = پهنای جدول
%\vspace{-10pt} % برای تنظیم فاصله عمودی (اختیاری)
%\begin{table}[h!]
\caption{داده‌های باردار شدن خازن \(C_1 = 20~\mu\mathrm{F}\) و نسبت \(V/V_0\) برحسبِ زمان}
\label{tab1}
\centering
\begin{tabular}{ccc}
\hline
$t \;[\mathrm{s}]$ & $V \;[\mathrm{V}]$ & $ V / V_0$ \\
\hline
$0$ & $10$ & $1.000$ \\
$15$ & $9.19$ & $0.919$ \\
$30$ & $8.46$ & $0.846$ \\
$45$ & $7.79$ & $0.779$ \\
$60$ & $7.17$ & $0.717$ \\
$75$ & $6.59$ & $0.659$ \\
$90$ & $6.07$ & $0.607$ \\
$105$ & $5.59$ & $0.559$ \\
$120$ & $5.14$ & $0.514$ \\
$135$ & $4.73$ & $0.473$ \\
\hline
\end{tabular}
%\end{table}
\end{wraptable}

مطابق شکل، خازنی با ظرفیتِ $C_1 = 20~\mu\mathrm{F}$ را به‌صورتِ سری به منبع تغذیهٔ ۱۰ ولتی و یک ولت‌متر وصل کردیم و برای زمان‌های مختلفِ جدول \ref{tab1}، عددی را که ولت‌متر نشان می‌دهد گزارش کردیم.


بین اختلاف پتانسیلِ منبع تغذیه $V_0$، اختلاف پتانسیلِ دو صفحهٔ خازن $V_C$ و عددی که ولت‌متر نشان می‌دهد $V$ رابطهٔ زیر برقرار است:
\[
V = V_0 - V_C = V_0 e^{-t/RC},
\]


\begin{figure}
\caption{{$\dfrac{V}{V_0}$ برحسب زمان}}
\label{fig1}
\centering
\pgfplotsset{
    width=12cm,
    height=5cm,
    scale only axis,
    tick align = outside,
    yticklabel style={/pgf/number format/fixed},
}
\begin{tikzpicture}
\begin{semilogyaxis}[
    %title=\rl{{$\dfrac{V}{V_0}$ برحسب زمان}},
    xlabel={\lr{[$\mathrm{s}$]} زمان $t$},
    ylabel={$\dfrac{V}{V_0}$},
    xmin=0, xmax=135,
    ymin=0.1, ymax=1,
    xtick distance=15,
    yminorgrids=true,
    ymajorgrids=true,
    xmajorgrids=true,
    grid style=dashed,
    grid=both,
    minor y tick num=9,
    minor x tick num=2,
    legend pos=north east,
    xticklabel style={
        /pgf/number format/fixed,
        /pgf/number format/precision=0,
        /pgf/number format/fixed zerofill
    },
]

\addplot[
    mark=square,
    smooth
]
coordinates {
    (0,   1.000)
    (15,  0.919)
    (30,  0.846)
    (45,  0.779)
    (60,  0.717)
    (75,  0.659)
    (90,  0.607)
    (105, 0.559)
    (120, 0.514)
    (135, 0.473)
};

%\legend{$V/V_0$}

\end{semilogyaxis}
\end{tikzpicture}
\end{figure}

در شکل \ref{fig1} منحنیِ تغییراتِ $\frac{V}{V_0}$ برحسبِ زمان در \emph{مقیاس‌بندیِ نیمه‌لگاریتمی} نمایش داده شده است.


حال با استفاده از شیبِ خط (با روشِ کمینهٔ مربعات) $\tau$، ثابت زمانیِ مدار را به‌دست می‌آوریم و مقاومت داخلیِ ولت‌متر را محاسبه می‌کنیم.

\subsection*{تعیینِ ثابتِ زمانیِ مدار و مقاومتِ داخلیِ ولت‌متر}

رابطهٔ تئوری مدار به صورت زیر است:
\begin{equation}
V = V_0 - V_C = V_0 e^{-t/RC},
\end{equation}
در نتیجه
\begin{equation}
\frac{V}{V_0} = e^{-t/RC}.
\end{equation}
اگر لگاریتم دهدهی بگیریم، داریم:
\begin{equation}
\log_{10}\!\left(\frac{V}{V_0}\right)
  = -\frac{t}{RC\ln 10}.
\end{equation}
با تعریف
\[
x_i = t_i, \qquad
y_i = \log_{10}\!\left(\frac{V_i}{V_0}\right),
\]
باید رابطهٔ خطی
\begin{equation}
y_i = a + b x_i
\end{equation}
صدق کند، که در آن
\begin{equation}
b = -\frac{1}{RC\ln 10} = -\frac{1}{\tau \ln 10}.
\end{equation}

برای به‌دست‌آوردن \(a\) و \(b\) از روش کمینه مربعات استفاده می‌کنیم. فرمول‌های شیب و عرض از مبدأ به صورت زیر است:
\begin{align}
b &=
\frac{n\sum x_i y_i - \left(\sum x_i\right)\left(\sum y_i\right)}
     {n\sum x_i^2 - \left(\sum x_i\right)^2},\\[2mm]
a &= \frac{1}{n}\left(\sum y_i - b\sum x_i\right),
\end{align}
که در آن \(n\) تعداد نقاط آزمایش است. با استفاده از داده‌های جدول \ref{tab1} داریم:
\begin{align}
\sum x_i &= \sum t_i = 675~\mathrm{s},\\
\sum x_i^2 &= \sum t_i^2 = 64125~\mathrm{s^2},\\
\sum y_i &\approx -1.6269,\\
\sum x_i y_i &\approx -154.48.
\end{align}
بنابراین:
\begin{align}
b &=
\frac{10(-154.48) - 675(-1.6269)}
     {10(64125) - 675^2} \\
&\approx \frac{-446.6}{185625} \\
&\approx -2.41\times10^{-3}\ \mathrm{s^{-1}},
\end{align}
و
\begin{align}
a &= \frac{1}{10}\left(\sum y_i - b\sum x_i\right) \\
&\approx -3.1\times10^{-4},
\end{align}
که مقدار \(a\) بسیار کوچک بوده و با مقدار تئوری صفر سازگار است. در نتیجه معادلهٔ خط برازش‌شده روی نمودار نیمه‌لگاریتمی به صورت زیر است:
\begin{equation}
\log_{10}\!\left(\frac{V}{V_0}\right)
  \approx -2.41\times10^{-3}\, t - 3.1\times10^{-4}.
\end{equation}

از رابطهٔ
\begin{equation}
b = -\frac{1}{\tau \ln 10}
\end{equation}
ثابت زمانی مدار به صورت زیر به‌دست می‌آید:
\begin{align}
\tau &= -\frac{1}{b\ln 10} \\
&= -\frac{1}{(-2.41\times10^{-3})\times 2.3026} \\
&\approx 1.8\times10^{2}\ \mathrm{s}.
\end{align}

در این آزمایش تنها مقاومت سری، مقاومت داخلی ولت‌متر \(R_\mathrm{v}\) است، بنابراین:
\begin{equation}
\tau = R_\mathrm{v} C_1
\qquad\Rightarrow\qquad
R_\mathrm{v} = \frac{\tau}{C_1}.
\end{equation}
با توجه به
\[
C_1 = 20~\mu\mathrm{F} = 20\times10^{-6}~\mathrm{F},
\]
خواهیم داشت:
\begin{align}
R_\mathrm{v}
&= \frac{1.8\times10^{2}}{20\times10^{-6}} \\
&\approx 9.0\times10^{6}\ \Omega \\
&\approx 9~\mathrm{M}\Omega.
\end{align}

در نتیجه:
\begin{equation}
\boxed{\tau \approx 1.8\times10^{2}\ \mathrm{s}}
\qquad
\boxed{R_\mathrm{v} \approx 9~\mathrm{M}\Omega}.
\end{equation}



\section{بی‌بار شدن خازن $C_2$ و تعیین مقاومت داخلی ولت‌متر}
\begin{wraptable}{l}{4cm}
%\begin{table}[h!]
\caption{داده‌های بی‌بار شدن خازن \(C_2 = 4~\mu\mathrm{F}\) و ولتاژ برحسبِ زمان}
\label{tab2}
\centering
\begin{tabular}{cc}
\hline
$t \;[\mathrm{s}]$ & $V \;[\mathrm{V}]$ \\
\hline
  $0$ & $10$ \\
 $15$ & $7.36$ \\
 $30$ & $5.37$ \\
 $45$ & $3.92$ \\
 $60$ & $2.87$ \\
 $75$ & $2.11$ \\
 $90$ & $1.54$ \\
$105$ & $1.13$ \\
$120$ & $0.83$ \\
$135$ & $0.60$ \\
\hline
\end{tabular}
%\end{table}
\end{wraptable}

ابتدا خازنی با ظرفیتِ $C_2 = 4~\mu\mathrm{F}$ را به‌وسیلهٔ یک منبع تغذیهٔ ۱۰ ولتی باردار کردیم؛ سپس خازن را از منبع تغذیه جدا کرده و به ولت‌متر وصل نمودیم تا از راهِ مقاومت داخلیِ ولت‌متر تخلیه شود. برای زمان‌های مختلفِ جدول \ref{tab2}، عددی را که ولت‌متر نشان می‌دهد گزارش کردیم.

\begin{figure}
\caption{{$\dfrac{V}{V_0}$ برحسب زمان برای خازن $C_2$}}
\label{fig2}
\centering
\pgfplotsset{
    width=12cm,
    height=10cm,
    scale only axis,
    tick align = outside,
    yticklabel style={/pgf/number format/fixed},
}
\begin{tikzpicture}
\begin{semilogyaxis}[
    %title=\rl{{$\dfrac{V}{V_0}$ برحسب زمان برای خازن $C_2$}},
    xlabel={\lr{[$\mathrm{s}$]} زمان $t$},
    ylabel={$\dfrac{V}{V_0}$},
    xmin=0, xmax=135,
    ymin=0.01, ymax=1.1,
    xtick distance=15,
    ymajorgrids=true,
    yminorgrids=true,
    xmajorgrids=true,
    grid=both,
    grid style=dashed,
    minor y tick num=9,
    minor x tick num=2,
    xticklabel style={
        /pgf/number format/fixed,
        /pgf/number format/precision=0,
        /pgf/number format/fixed zerofill
    },
]

\addplot[
    mark=square,
    smooth
]
coordinates {
    (0,   1.000)
    (15,  0.736)
    (30,  0.537)
    (45,  0.392)
    (60,  0.287)
    (75,  0.211)
    (90,  0.154)
    (105, 0.113)
    (120, 0.083)
    (135, 0.060)
};

%\legend{$V/V_0$}

\end{semilogyaxis}
\end{tikzpicture}
\end{figure}


در شکل \ref{fig2} منحنیِ تغییراتِ $\frac{V}{V_0}$ برحسبِ زمان در \emph{مقیاس‌بندیِ نیمه‌لگاریتمی} نمایش داده شده است.


حال با استفاده از شیبِ خط (با روشِ کمینهٔ مربعات) $\tau$، ثابت زمانیِ مدار را به‌دست می‌آوریم و مقاومت داخلیِ ولت‌متر را محاسبه می‌کنیم.


\subsection*{تعیینِ ثابتِ زمانیِ مدار و مقاومتِ داخلیِ ولت‌متر (خازن \texorpdfstring{$C_2 = 4\,\mu\mathrm{F}$}{C2})}

در این آزمایش، خازن $C_2$ ابتدا تا ولتاژ $V_0 = 10~\mathrm{V}$ باردار شد و سپس از طریق مقاومت داخلی ولت‌متر تخلیه گردید. رابطهٔ تئوری ولتاژ روی خازن در هنگام تخلیه به صورت زیر است:
\begin{equation}
V(t) = V_0 e^{-t/RC},
\end{equation}
بنابراین
\begin{equation}
\frac{V}{V_0} = e^{-t/RC}.
\end{equation}
با گرفتن لگاریتم دهدهی، خواهیم داشت:
\begin{equation}
\log_{10}\!\left(\frac{V}{V_0}\right)
  = -\frac{t}{RC\ln 10}.
\end{equation}
اگر تعریف کنیم
\[
x_i = t_i, \qquad
y_i = \log_{10}\!\left(\frac{V_i}{V_0}\right),
\]
آنگاه داده‌ها باید روی یک خط راست از فرم
\begin{equation}
y_i = a + b x_i
\end{equation}
قرار گیرند که در آن
\begin{equation}
b = -\frac{1}{RC\ln 10} = -\frac{1}{\tau \ln 10}.
\end{equation}

\paragraph{ساختن داده‌های لگاریتمی}
با توجه به $V_0 = 10~\mathrm{V}$، نسبت‌های $\frac{V_i}{V_0}$ و لگاریتم دهدهی آن‌ها به صورت زیر به‌دست می‌آیند (چند رقم اعشار نگه داشته شده است):
\[
\begin{array}{c c c c}
\hline
t_i~[\mathrm{s}] & V_i~[\mathrm{V}] & V_i/V_0 & y_i=\log_{10}(V_i/V_0) \\
\hline
0   & 10.00 & 1.000 &  0.0000 \\
15  &  7.36 & 0.736 & -0.1331 \\
30  &  5.37 & 0.537 & -0.2700 \\
45  &  3.92 & 0.392 & -0.4067 \\
60  &  2.87 & 0.287 & -0.5421 \\
75  &  2.11 & 0.211 & -0.6757 \\
90  &  1.54 & 0.154 & -0.8125 \\
105 &  1.13 & 0.113 & -0.9469 \\
120 &  0.83 & 0.083 & -1.0809\\
135 &  0.60 & 0.060 & -1.2218\\
\hline
\end{array}
\]

\paragraph{روش کمینه مربعات}

برای خط $y = a + bx$، شیب $b$ و عرض از مبدأ $a$ با روش کمینه مربعات برابرند با:
\begin{align}
b &=
\frac{n\sum x_i y_i - \left(\sum x_i\right)\left(\sum y_i\right)}
     {n\sum x_i^2 - \left(\sum x_i\right)^2},\\[2mm]
a &= \frac{1}{n}\left(\sum y_i - b\sum x_i\right),
\end{align}
که در آن $n=10$ تعداد نقاط است.

با استفاده از داده‌ها داریم:
\begin{align}
\sum x_i &= \sum t_i = 675~\mathrm{s},\\
\sum x_i^2 &= \sum t_i^2 = 64125~\mathrm{s^2},\\
\sum y_i &\approx -6.090,\\
\sum x_i y_i &\approx -578.82.
\end{align}

در نتیجه شیب خط برابر است با:
\begin{align}
b &=
\frac{10(-578.82) - 675(-6.090)}
     {10(64125) - 675^2} \\[1mm]
&=
\frac{-5788.2 + 4109.0}{641250 - 455625} \\[1mm]
&\approx \frac{-1679.2}{185625} \\[1mm]
&\approx -9.04\times 10^{-3}\ \mathrm{s^{-1}}.
\end{align}

عرض از مبدأ نیز برابر است با:
\begin{align}
a &= \frac{1}{10}\left(\sum y_i - b\sum x_i\right)\\[1mm]
&\approx \frac{1}{10}\left(-6.090 - (-9.04\times10^{-3})\times 675\right)\\[1mm]
&\approx 1.0\times10^{-3},
\end{align}
که مقدار بسیار کوچکی است و با مقدار تئوری (تقریباً صفر) سازگار می‌باشد.

پس معادلهٔ خط برازش‌شده روی نمودار نیمه‌لگاریتمی به صورت زیر است:
\begin{equation}
\log_{10}\!\left(\frac{V}{V_0}\right)
  \approx -9.04\times10^{-3}\, t + 1.0\times10^{-3}.
\end{equation}

\paragraph{محاسبهٔ ثابت زمانی \texorpdfstring{$\tau$}{tau}}

از رابطهٔ
\begin{equation}
b = -\frac{1}{\tau \ln 10}
\end{equation}
خواهیم داشت:
\begin{align}
\tau &= -\frac{1}{b\ln 10} \\[1mm]
&= -\frac{1}{(-9.04\times10^{-3})\times 2.3026} \\[1mm]
&\approx 4.8\times10^{1}\ \mathrm{s} \\[1mm]
&\approx 48\ \mathrm{s}.
\end{align}

\paragraph{محاسبهٔ مقاومت داخلی ولت‌متر}

در این مدار تنها مقاومت موجود، مقاومت داخلی ولت‌متر $R_\mathrm{v}$ است، بنابراین:
\begin{equation}
\tau = R_\mathrm{v} C_2
\qquad\Rightarrow\qquad
R_\mathrm{v} = \frac{\tau}{C_2}.
\end{equation}
با توجه به
\[
C_2 = 4~\mu\mathrm{F} = 4\times10^{-6}~\mathrm{F},
\]
خواهیم داشت:
\begin{align}
R_\mathrm{v}
&= \frac{4.8\times10^{1}}{4\times10^{-6}} \\[1mm]
&\approx 1.2\times10^{7}\ \Omega \\[1mm]
&\approx 12~\mathrm{M}\Omega.
\end{align}

در نتیجه مقادیر به‌دست‌آمده برای این آزمایش به صورت زیر است:
\begin{equation}
\boxed{\tau \approx 4.8\times10^{1}\ \mathrm{s}}
\qquad
\boxed{R_\mathrm{v} \approx 12~\mathrm{M}\Omega}.
\end{equation}



حال مقدار متوسط مقاومت داخلی ولت‌متر $(\bar{R})$ را با استفاده از نتایج به‌دست آمده از قسمت‌های قبلی، تعیین می‌کنیم.

\subsection*{تعیین مقدار متوسط مقاومت داخلی ولت‌متر}

در دو بخشِ قبلی، با استفاده از دو خازنِ مختلف، مقدارِ مقاومتِ داخلیِ ولت‌متر به‌صورتِ زیر به‌دست آمد:
\[
R_1 \approx 9~\mathrm{M\Omega}, \qquad
R_2 \approx 12~\mathrm{M\Omega}.
\]
برای به‌دست‌آوردن مقدار متوسط مقاومت داخلی ولت‌متر، از رابطهٔ
\[
\bar{R} = \frac{R_1 + R_2}{2}
\]
استفاده می‌کنیم. بنابراین:
\[
\bar{R} = \frac{9 + 12}{2}~\mathrm{M\Omega}
       = 10.5~\mathrm{M\Omega}.
\]
در نتیجه مقدار متوسط مقاومت داخلی ولت‌متر برابر است با:
\[
\boxed{\bar{R} = 10.5~\mathrm{M\Omega}}.
\]


\section{بررسیِ تجربیِ ظرفیتِ معادلِ خازن‌های سری}
\begin{wraptable}{l}{4cm}

%\begin{table}[h!]
\caption{داده‌های تخلیهٔ هم‌زمان خازن‌های \(C_1\) و \(C_2\) در اتصالِ سری و ولتاژ برحسبِ زمان}
\label{tab3}
\centering
\begin{tabular}{cc}
\hline
$t \;[\mathrm{s}]$ & $V \;[\mathrm{V}]$ \\
\hline
  $0$ & $10$ \\
 $15$ & $6.95$ \\
 $30$ & $4.63$ \\
 $45$ & $3.12$ \\
 $60$ & $2.11$ \\
 $75$ & $1.42$ \\
 $90$ & $0.96$ \\
$105$ & $0.65$ \\
$120$ & $0.48$ \\
$135$ & $0.30$ \\
\hline
\end{tabular}
%\end{table}
\end{wraptable}

خازن‌های $C_1$ و $C_2$ که در قسمت‌های پیشین مورد استفاده قرار دادیم را به‌صورتِ سری به یک منبع تغذیهٔ ۱۰ ولتی و یک ولت‌متر وصل کردیم و برای زمان‌های مختلفِ جدول \ref{tab3}، عددِ نمایش‌داده‌شدهٔ ولت‌متر را گزارش کردیم.


حال منحنیِ تغییراتِ $V$ برحسبِ زمان را رسم می‌کنیم؛ سپس با استفاده از منحنی و تعریفِ $\tau$، ثابت زمانیِ مدار را به‌دست می‌آوریم.

\subsection*{تعیینِ ثابتِ زمانیِ مدار از روی منحنی $V(t)$}

در مدار تخلیهٔ خازن، ولتاژ دو سر خازن بر حسب زمان از رابطهٔ
\begin{equation}
V(t) = V_0 e^{-t/\tau}
\end{equation}
پیروی می‌کند که در آن $\tau = RC$ ثابت زمانی مدار است. طبق این رابطه، در لحظهٔ
\begin{equation}
t = \tau \quad \Rightarrow \quad V(\tau) = \frac{V_0}{e}.
\end{equation}

در این آزمایش $V_0 = 10~\mathrm{V}$ است؛ بنابراین:
\begin{equation}
V(\tau) = \frac{10}{e} \approx 3.68~\mathrm{V}.
\end{equation}

\begin{figure}[h!]
\caption{{$V(t)$ برای تخلیهٔ خازن‌های $C_1$ و $C_2$ (اتصال سری)}},
\label{fig3}
\centering
\pgfplotsset{
    width=12cm,
    height=7cm,
    scale only axis,
    tick align = outside,
    yticklabel style={/pgf/number format/fixed},
}
\begin{tikzpicture}
\begin{axis}[
    xlabel={\lr{[$\mathrm{s}$]} زمان $t$},
    ylabel={\lr{[$\mathrm{V}$]} ولتاژ $V$},
    xmin=0, xmax=135,
    ymin=0, ymax=10.5,
    xtick distance=15,
    ytick distance=1,
    ymajorgrids=true,
    xmajorgrids=true,
    minor y tick num=1,
    minor x tick num=1,
    grid style=dashed,
]

\addplot[
    mark=square,
    only marks
]
coordinates {
    (0,   10.00)
    (15,  6.95)
    (30,  4.63)
    (45,  3.12)
    (60,  2.11)
    (75,  1.42)
    (90,  0.96)
    (105, 0.65)
    (120, 0.48)
    (135, 0.30)
};

\addplot[
    smooth,
    domain=0:135,
    samples=100
]
{10.03282*exp(x / -38.66179) };

\legend{\rl{داده‌های تجربی},
$y=10.03282 \exp{\frac{x}{-38.66179}}$
}

\end{axis}
\end{tikzpicture}
\end{figure}


با توجه به جدول \ref{tab3} و منحنی $V(t)$ رسم‌شده (شکل \ref{fig3})، مشاهده می‌شود که این مقدار ولتاژ بین دو نقطهٔ زیر قرار می‌گیرد:
\[
\begin{aligned}
t_1 &= 30~\mathrm{s}, & V_1 &= 4.63~\mathrm{V},\\
t_2 &= 45~\mathrm{s}, & V_2 &= 3.12~\mathrm{V}.
\end{aligned}
\]

برای تخمین زمان $\tau$، بین این دو نقطه درون‌یابی خطی انجام می‌دهیم. ابتدا شیب خط بین دو نقطه را حساب می‌کنیم:
\begin{equation}
m = \frac{V_2 - V_1}{t_2 - t_1}
  = \frac{3.12 - 4.63}{45 - 30}
  = \frac{-1.51}{15}
  \approx -0.1007~\mathrm{V/s}.
\end{equation}

اختلاف ولتاژ از مقدار $V_1$ تا $V(\tau) = 3.68~\mathrm{V}$ برابر است با:
\begin{equation}
\Delta V = V(\tau) - V_1 = 3.68 - 4.63 = -0.95~\mathrm{V}.
\end{equation}
پس زمان لازم برای این تغییر ولتاژ:
\begin{equation}
\Delta t = \frac{\Delta V}{m}
         = \frac{-0.95}{-0.1007}
         \approx 9.45~\mathrm{s}.
\end{equation}

در نتیجه زمان متناظر با $V = V_0/e$ برابر است با:
\begin{equation}
\tau \approx t_1 + \Delta t = 30~\mathrm{s} + 9.45~\mathrm{s}
     \approx 3.9\times 10^{1}~\mathrm{s}.
\end{equation}

بنابراین، ثابت زمانیِ مدار به‌صورتِ زیر به‌دست می‌آید:
\begin{equation}
\boxed{\tau \approx 3.9\times 10^{1}\ \mathrm{s} \; (\approx 39~\mathrm{s}).}
\end{equation}



حال با استفاده از $\tau$ِ به‌دست‌آمده و مقاومت داخلیِ ولت‌متر، ظرفیتِ معادلِ خازن‌ها را به‌دست می‌آوریم و خطای این ظرفیت را نسبت به ظرفیتِ معادلِ 
\(\frac{1}{C} = \frac{1}{C_1} + \frac{1}{C_2}\)
محاسبه می‌کنیم.

\subsection*{تعیینِ ظرفیتِ معادلِ خازن‌ها و محاسبهٔ خطا}

از آزمایش قبلی (تخلیهٔ هم‌زمان دو خازن) ثابت زمانی مدار به صورت
\[
\tau \approx 3.9\times 10^{1}~\mathrm{s}
\]
به‌دست آمد. همچنین در دو قسمت قبلی میانگین مقاومت داخلی ولت‌متر را
\[
\bar{R} \approx 10.5\times 10^{6}~\Omega
\]
به‌دست آوردیم. چون در این مدار تنها مقاومت سری، مقاومت داخلی ولت‌متر است، طبق تعریف ثابت زمانی داریم:
\begin{equation}
\tau = \bar{R}\,C_{\mathrm{eq}} \qquad\Rightarrow\qquad
C_{\mathrm{eq}}^{(\mathrm{exp})} = \frac{\tau}{\bar{R}}.
\end{equation}
با جایگذاری مقادیر عددی:
\begin{align}
C_{\mathrm{eq}}^{(\mathrm{exp})}
&= \frac{3.9\times 10^{1}}{10.5\times 10^{6}}~\mathrm{F} \\[1mm]
&= \frac{3.9}{10.5}\times 10^{-5}~\mathrm{F} \\[1mm]
&\approx 3.71\times 10^{-6}~\mathrm{F} \\[1mm]
&\approx 3.71~\mu\mathrm{F}.
\end{align}

ظرفیت معادل تئوریِ دو خازن مطابق رابطهٔ
\begin{equation}
\frac{1}{C} = \frac{1}{C_1} + \frac{1}{C_2}
\end{equation}
به‌صورت زیر است. با در نظر گرفتن
\[
C_1 = 20~\mu\mathrm{F}, \qquad C_2 = 4~\mu\mathrm{F},
\]
داریم:
\begin{align}
\frac{1}{C_{\mathrm{eq}}^{(\mathrm{th})}}
&= \frac{1}{20} + \frac{1}{4} \quad (\mu\mathrm{F}^{-1}) \\[1mm]
&= 0.05 + 0.25 = 0.30, \\[1mm]
\Rightarrow\quad
C_{\mathrm{eq}}^{(\mathrm{th})}
&= \frac{1}{0.30} \approx 3.33~\mu\mathrm{F}.
\end{align}

اکنون خطای نسبی ظرفیت معادل به‌دست‌آمده از آزمایش نسبت به مقدار تئوری را محاسبه می‌کنیم:
\begin{equation}
\delta C = \frac{\bigl|C_{\mathrm{eq}}^{(\mathrm{exp})}
             - C_{\mathrm{eq}}^{(\mathrm{th})}\bigr|}
            {C_{\mathrm{eq}}^{(\mathrm{th})}} \times 100~\%.
\end{equation}
با جایگذاری:
\begin{align}
\delta C
&= \frac{|3.71 - 3.33|}{3.33}\times 100~\% \\[1mm]
&\approx \frac{0.38}{3.33}\times 100~\% \\[1mm]
&\approx 1.14\times 10^{1}~\% \\[1mm]
&\approx 11~\%.
\end{align}

در نتیجه:
\begin{equation}
\boxed{C_{\mathrm{eq}}^{(\mathrm{exp})} \approx 3.71~\mu\mathrm{F}}
\qquad
\boxed{C_{\mathrm{eq}}^{(\mathrm{th})} \approx 3.33~\mu\mathrm{F}}
\qquad
\boxed{\delta C \approx 11~\%}.
\end{equation}



\section{بررسیِ تجربیِ ظرفیتِ معادلِ خازن‌های موازی}

\begin{wraptable}{l}{4cm}

%\begin{table}[h!]
\caption{داده‌های تخلیهٔ خازن‌های \(C_1\) و \(C_2\) در اتصالِ موازی و ولتاژ برحسبِ زمان}
\label{tab4}
\centering
\begin{tabular}{cc}
\hline
$t \;[\mathrm{s}]$ & $V \;[\mathrm{V}]$ \\
\hline
  $0$ & $9.83$ \\
 $30$ & $8.65$ \\
 $60$ & $7.59$ \\
 $90$ & $6.65$ \\
$120$ & $5.83$ \\
$150$ & $5.11$ \\
$180$ & $4.48$ \\
$210$ & $3.93$ \\
$240$ & $3.46$ \\
$270$ & $3.02$ \\
$300$ & $2.65$ \\
\hline
\end{tabular}
%\end{table}
\end{wraptable}

خازن‌های $C_1$ و $C_2$ را به‌صورتِ موازی به‌هم بسته و سپس به‌وسیلهٔ یک منبع تغذیهٔ ۱۰ ولتی باردار می‌کنیم. خازن‌ها را از منبع جدا کرده و به ولت‌متر وصل می‌کنیم تا خازن‌ها از راهِ مقاومت داخلیِ ولت‌متر، تخلیه شوند و برای زمان‌های مختلفِ جدول \ref{tab4}، عددی را که ولت‌متر نمایش می‌دهد گزارش کردیم.



شکل \ref{fig4}، منحنیِ نمایشِ تغییراتِ $V$ بر حسبِ زمان را نشان می‌دهد. با استفاده از این منحنی و تعریفِ $\tau$، ثابتِ زمانی را به‌دست می‌آوریم.


\subsection*{تعیینِ ثابتِ زمانیِ مدار از روی منحنی $V(t)$ (خازن‌های $C_1$ و $C_2$ موازی)}

در تخلیهٔ خازن، ولتاژ دو سر خازن معادل بر حسب زمان از رابطهٔ
\begin{equation}
V(t) = V_0 e^{-t/\tau}
\end{equation}
پیروی می‌کند که در آن $\tau$ ثابت زمانی مدار است. طبق این رابطه،
\begin{equation}
V(\tau) = \frac{V_0}{e}.
\end{equation}
در این آزمایش، از دادهٔ اولیه در زمان $t=0$ داریم:
\[
V_0 \approx 9.83~\mathrm{V},
\]
بنابراین:
\begin{equation}
V(\tau) = \frac{9.83}{e} \approx 3.62~\mathrm{V}.
\end{equation}

\begin{figure}[h!]
\caption{{$V(t)$ برای تخلیهٔ خازن‌های $C_1$ و $C_2$ (اتصال موازی)}},
\label{fig4}
\centering
\pgfplotsset{
    width=12cm,
    height=7cm,
    scale only axis,
    tick align = outside,
    yticklabel style={/pgf/number format/fixed},
}
\begin{tikzpicture}
\begin{axis}[
    xlabel={\lr{[$\mathrm{s}$]} زمان $t$},
    ylabel={\lr{[$\mathrm{V}$]} ولتاژ $V$},
    xmin=0, xmax=300,
    ymin=0, ymax=10.5,
    xtick distance=30,
    ytick distance=1,
    ymajorgrids=true,
    xmajorgrids=true,
    minor y tick num=1,
    minor x tick num=1,
    grid style=dashed,
]

\addplot[
    mark=square,
    only marks
]
coordinates {
    (0, 9.83)
    (30, 8.65)
    (60, 7.59)
    (90, 6.65)
    (120, 5.83)
    (150, 5.11)
    (180, 4.48)
    (210, 3.93)
    (240, 3.46)
    (270, 3.02)
    (300, 2.65)
};

\addplot[
    smooth,
    domain=0:300,
    samples=100
]
{9.84072*exp(x / -228.9636) };

\legend{\rl{داده‌های تجربی}, 
$y=9.84072 \exp{\frac{x}{-228.9636}}$
}

\end{axis}
\end{tikzpicture}
\end{figure}

با توجه به جدول \ref{tab4} و منحنی $V(t) $(شکل \ref{fig4}) مشاهده می‌شود که این مقدار ولتاژ بین دو نقطهٔ
\[
\begin{aligned}
t_1 &= 210~\mathrm{s}, & V_1 &= 3.93~\mathrm{V},\\
t_2 &= 240~\mathrm{s}, & V_2 &= 3.46~\mathrm{V}
\end{aligned}
\]
قرار دارد. برای تخمین زمان $\tau$ بین این دو نقطه درون‌یابی خطی انجام می‌دهیم. شیب خط بین دو نقطه برابر است با:
\begin{equation}
m = \frac{V_2 - V_1}{t_2 - t_1}
  = \frac{3.46 - 3.93}{240 - 210}
  = \frac{-0.47}{30}
  \approx -1.57\times 10^{-2}~\mathrm{V/s}.
\end{equation}
اختلاف ولتاژ از $V_1$ تا $V(\tau)$ برابر است با:
\begin{equation}
\Delta V = V(\tau) - V_1 = 3.62 - 3.93 \approx -0.31~\mathrm{V}.
\end{equation}
پس زمان لازم برای این تغییر ولتاژ:
\begin{equation}
\Delta t = \frac{\Delta V}{m}
         \approx \frac{-0.31}{-1.57\times 10^{-2}}
         \approx 2.0\times 10^{1}~\mathrm{s}.
\end{equation}
در نتیجه زمان متناظر با $V = V_0/e$ (یعنی ثابت زمانی مدار) برابر است با:
\begin{equation}
\tau \approx t_1 + \Delta t
     \approx 210~\mathrm{s} + 20~\mathrm{s}
     \approx 2.3\times 10^{2}~\mathrm{s}.
\end{equation}

بنابراین ثابت زمانی مدار خازن‌های $C_1$ و $C_2$ (متصل به صورت موازی) به صورت زیر به‌دست می‌آید:
\begin{equation}
\boxed{\tau \approx 2.3\times 10^{2}\ \mathrm{s} \; (\approx 230~\mathrm{s}).}
\end{equation}


اکنون با استفاده از $\tau$ و مقاومت داخلیِ ولت‌متر، ظرفیتِ معادلِ خازن‌ها را به‌دست آورده و خطای آن را با ظرفیتِ معادلِ $C = C_1 + C_2$ محاسبه می‌کنیم.

\subsection*{تعیین ظرفیت معادل خازن‌ها از روی $\tau$ و مقایسه با حالت موازی}

در این بخش، خازن‌های $C_1$ و $C_2$ به صورت موازی بسته شده و از طریق مقاومت داخلی ولت‌متر تخلیه شده‌اند. از بخش قبل، ثابت زمانی مدار را
\[
\tau \approx 2.3\times 10^{2}~\mathrm{s}
\]
به‌دست آوردیم. همچنین میانگین مقاومت داخلی ولت‌متر از دو آزمایش قبلی برابر بود با
\[
\bar{R} \approx 10.5\times 10^{6}~\Omega.
\]

چون در مدار فقط همین مقاومت سری وجود دارد، تعریف ثابت زمانی می‌دهد:
\begin{equation}
\tau = \bar{R}\,C_{\mathrm{eq}}
\qquad\Rightarrow\qquad
C_{\mathrm{eq}}^{(\mathrm{exp})} = \frac{\tau}{\bar{R}}.
\end{equation}
با جایگذاری مقادیر عددی:
\begin{align}
C_{\mathrm{eq}}^{(\mathrm{exp})}
&= \frac{2.3\times 10^{2}}{10.5\times 10^{6}}~\mathrm{F} \\[1mm]
&= \frac{2.3}{10.5}\times 10^{-4}~\mathrm{F} \\[1mm]
&\approx 2.19\times 10^{-5}~\mathrm{F} \\[1mm]
&\approx 21.9~\mu\mathrm{F}.
\end{align}

از طرف دیگر، ظرفیت معادل تئوری دو خازن موازی برابر است با:
\begin{equation}
C_{\mathrm{eq}}^{(\mathrm{th})} = C_1 + C_2.
\end{equation}
با در نظر گرفتن
\[
C_1 = 20~\mu\mathrm{F}, \qquad C_2 = 4~\mu\mathrm{F},
\]
خواهیم داشت:
\begin{equation}
C_{\mathrm{eq}}^{(\mathrm{th})} = 20 + 4 = 24~\mu\mathrm{F}.
\end{equation}

اکنون خطای نسبی ظرفیت معادل اندازه‌گیری‌شده نسبت به مقدار تئوری را محاسبه می‌کنیم:
\begin{equation}
\delta C = \frac{\bigl|C_{\mathrm{eq}}^{(\mathrm{exp})}
             - C_{\mathrm{eq}}^{(\mathrm{th})}\bigr|}
            {C_{\mathrm{eq}}^{(\mathrm{th})}} \times 100~\%.
\end{equation}
با جایگذاری:
\begin{align}
\delta C
&= \frac{|21.9 - 24.0|}{24.0}\times 100~\% \\[1mm]
&\approx \frac{2.1}{24.0}\times 100~\% \\[1mm]
&\approx 8.7~\%.
\end{align}

در نتیجه:
\begin{equation}
\boxed{C_{\mathrm{eq}}^{(\mathrm{exp})} \approx 21.9~\mu\mathrm{F}}
\qquad
\boxed{C_{\mathrm{eq}}^{(\mathrm{th})} = 24~\mu\mathrm{F}}
\qquad
\boxed{\delta C \approx 8.7~\%.}
\end{equation}

\paragraph{یادداشت}
نمودار‌های رسم در شکل های \ref{fig3} و \ref{fig4} ضابطه‌شان با وارد کردن داده‌های جداول، در یک اسکریپت پایتون \footnote{کد پایتون در آدرس \lr{\url{https://abulfadl.ir/genphyslabs2/4/LMA}} قابل مشاهده است.} به‌دست آمده‌اند. در این برنامه از \lr{ Levenberg–Marquardt algorithm} برای این منظور استفاده می‌کند.

مانند سایر الگوریتم‌های کمینه‌سازی عددی، الگوریتم لونبرگ-مارکارد یک رویهٔ تکراری است. برای شروع کمینه‌سازی، کاربر باید یک حدس آغازین برای بردار پارامترها
$\mathbf{p}$ ارائه کند. در بسیاری مواقع یک حدس ناآگاهانهٔ استاندارد مانند
\[
\mathbf{p}^{T} = (1,1,\dots,1)
\]
به‌خوبی کار می‌کند. در جاهای دیگر، الگوریتم تنها وقتی کار می‌کند که حدس آغازین تا حدی به جواب نهایی نزدیک باشد.

در هر گام تکرار، بردار پارامترها $\mathbf{p}$ با یک تخمین جدید $\mathbf{p+q}$ جایگزین می‌شود. برای دستیابی به
$\mathbf{q}$، توابع $f_i(\mathbf{p+q})$ با خطی‌سازی‌شان
\[
\mathbf{f(p+q)} \approx \mathbf{f(p)} + J \mathbf{q}
\]
تخمین زده می‌شوند که $J$ ژاکوبینِ $\mathbf{f}$ در $\mathbf{p}$ است.

در یک کمینهٔ مجموع مربعات $S$، داریم
\[
\nabla_{\mathbf{q}} S = 0.
\]
با خطی‌سازی بالا، از این به معادلهٔ زیر می‌رسیم:
\[
(J^{T} J)\,\mathbf{q} = - J^{T} \mathbf{f}
\]
که $\mathbf{q}$ را می‌توان از آن با وارون کردن $J^{T} J$ به‌دست آورد.
کلید LMA جایگزینی این معادله با «نسخهٔ میراشده»ٔ آن است:
\[
(J^{T} J + \lambda)\,\mathbf{q} = - J^{T} \mathbf{f}
\]
ضریب میرایی نامنفی $\lambda$ در هر تکرار تنظیم می‌شود. اگر کاهش $S$
تند بود، مقدار کوچک‌تری به آن می‌دهیم که الگوریتم را به GNA نزدیک‌تر می‌کند،
اما اگر یک تکرار کاهش ناکافی نشان داد، $\lambda$ را افزایش داده و یک گام را
نزدیک‌تر به جهت نزول گرادیانی برمی‌داریم. ضریب میرایی مشابهی در ساماندهی
تیخونوف دیده می‌شود که در حل مسائل خطی کژنمود سودمند است.

اگر یک طول قدم بازیابی‌شده یا کاهش مجموع مربعات برای آخرین مجموعهٔ پارامترهای
$\mathbf{p}$ از مقادیر از پیش تعیین‌شده کمتر باشند، تکرار پایان می‌یابد و آخرین
بردار پارامتر $\mathbf{p}$ به عنوان جواب در نظر گرفته می‌شود.



\end{document}